\documentclass[UTF8]{ctexart}
\usepackage{xeCJK}
\usepackage{geometry}
\geometry{left=2cm,right=2cm,top=2.5cm,bottom=2.5cm}
\setlength{\headheight}{12.7pt}

\usepackage{titlesec}
\titleformat{\section}
  {\normalfont\large\bfseries}
  {\thesection}
  {1em}
  {}
\titlespacing*{\section}
  {0pt}
  {3.5ex plus 1ex minus .2ex}
  {2.3ex plus .2ex}

\usepackage{amsmath}
\usepackage{amssymb}
\usepackage{amsthm}
\usepackage{enumitem}
\setlist[enumerate]{itemsep=0pt,topsep=0pt,parsep=0pt,leftmargin=0pt,itemindent=2.5em}
\setlist[itemize]{itemsep=0pt,topsep=0pt,parsep=0pt,leftmargin=0pt,itemindent=2.5em}
\usepackage{fancyhdr}
\pagestyle{fancy}
\fancyhf{}
\usepackage{graphicx}
\usepackage{hyperref}
\usepackage{indentfirst}
\usepackage{lmodern}


\begin{document}
\setlength{\abovedisplayskip}{1pt}
\setlength{\belowdisplayskip}{1pt}

\section{序数乘法}

\hypertarget{sec:定义1.1}{}
\noindent\textbf{定义1.1 (序数乘法)}

给定\href{../02 序数/01 序数#nameddest=sec:定义1.1}{序数} $\alpha$,
递归定义类映射$\mathbf{Mul}_\alpha:\mathbf{On}\to\mathbf{On}$如下:
\vspace{3pt}
\[
    \forall\beta\in\mathbf{On}:\quad
    \mathbf{Mul}_\alpha(\beta)\overset{\mathrm{def}}{=}\left\{\begin{aligned}
        &0,\quad&&\beta=0,\\[-3pt]
        &\mathbf{Mul}_\alpha(\max\beta)+\alpha,\quad&&
        \beta\text{\,\,是后继序数},\\[-3pt]
        &\sup\mathbf{Mul}_\alpha[\beta],\quad&&
        \beta\text{\,\,是极限序数}.
    \end{aligned}\right.
\]

接下来, 对任意的序数$\alpha,\beta$, 定义
\[
    \alpha\cdot\beta\overset{\mathrm{def}}{=}\mathbf{Mul}_\alpha(\beta),
\]
也可写作$\alpha\beta$, 称为\textbf{序数乘法}.

\vspace{10pt}

\hypertarget{sec:定理1.1}{}
\noindent\textbf{定理1.1 (序数乘法)}

任给序数$\alpha,\beta$,
\begin{enumerate}
    \item $\alpha0=0$, $\alpha1=\alpha$,
    \item $\alpha\beta^+=\alpha\beta+\alpha$,
    \item 如果$\beta$是极限序数,
    那么$\displaystyle\alpha\beta=\sup_{\gamma<\beta}\alpha\gamma$.
\end{enumerate}

\vspace{15pt}

显然\href{../../03 自然数/02 自然数算术/02 自然数的乘法#nameddest=sec:定义1.1}
{自然数的乘法}是序数乘法的特例. 易证任给序数$\alpha$有$0\alpha=0,1\alpha=\alpha$.

\vspace{10pt}

\hypertarget{sec:定理1.2}{}
\noindent\textbf{定理1.2}

任给序数$\alpha,\beta,\beta'$, 如果$\alpha\neq 0$且$\beta<\beta'$,
那么$\alpha\beta<\alpha\beta'$.

\noindent{\kaishu 证明.}

给定$\alpha\neq 0$和$\beta$, 对$\beta'>\beta$归纳证明$\alpha\beta<\alpha\beta'$.
任取$\beta'>\beta$, 假设
\[
    \forall\gamma:\quad\beta<\gamma<\beta'\to\alpha\beta<\alpha\gamma.
\]

\begin{enumerate}
    \item 如果$\beta'$是后继序数, 即有$\beta'=\gamma^+$,
    那么有$\beta\leqslant\gamma<\beta'$, 依假设有
    \[
        \alpha\beta\leqslant\alpha\gamma<\alpha\gamma^+=\alpha\beta'.
    \]
    \item 如果$\beta'$是极限序数, 则$\beta^+<\beta'$, 从而
    \begin{equation*}
        \alpha\beta<\alpha\beta^+\leqslant
        \sup_{\gamma<\beta'}\alpha\gamma=\alpha\beta'.
    \tag*{\qedsymbol}\end{equation*}
\end{enumerate}

\vspace{10pt}

\hypertarget{sec:定理1.3}{}
\noindent\textbf{定理1.3}

任给序数$\alpha,\alpha',\beta$, 如果$\alpha<\alpha'$,
那么$\alpha\beta\leqslant\alpha'\beta$.

\noindent{\kaishu 证明.}

给定$\alpha<\alpha'$, 对$\beta$归纳证明$\alpha\beta\leqslant\alpha'\beta$.
任取$\beta$, 假设对任意的$\gamma<\beta$有$\alpha\gamma\leqslant\alpha'\gamma$.

\begin{enumerate}
    \item 如果$\beta=0$, 显然.
    \item 如果$\beta$是后继序数, 即有$\beta=\gamma^+$, 那么$\gamma<\beta$, 依假设有
    \[
        \alpha\beta=\alpha\gamma+\alpha\leqslant\alpha'\gamma+\alpha'=\alpha'\beta.
        \\[1pt]
    \]
    \item 如果$\beta$是极限序数, 那么依假设有
    \vspace{-1pt}
    \begin{equation*}
        \alpha\beta=\sup_{\gamma<\beta}\alpha\gamma\leqslant
        \sup_{\gamma<\beta}\alpha'\gamma=\alpha'\beta.
    \tag*{\qedsymbol}\end{equation*}
\end{enumerate}





\newpage

\hypertarget{sec:定理1.4}{}
\noindent\textbf{定理1.4 (序数带余除法)}

任给序数$\alpha,\beta$, 如果$\beta\neq 0$, 那么存在唯一的一对序数$\eta,\varepsilon$
使得
\[
    \alpha=\beta\eta+\varepsilon\quad\land\quad\varepsilon<\beta,
\]
分别记作$\alpha/\beta$和$\alpha\%\beta$.

\noindent{\kaishu 证明.}

\noindent{\kaishu 存在性.}

取满足$\beta\theta>\alpha$的最小的$\theta$,
即对任意的$\eta<\theta$有$\beta\eta\leqslant\alpha$. 显然$\theta$非零.

\vspace{3pt}
如果$\theta$是极限序数, 那么$\displaystyle\beta\theta=\sup_{\eta<\theta}\beta\eta
\leqslant\alpha$矛盾. 所以$\theta$是后继序数, 有$\theta=\eta^+$.

\vspace{3pt}
于是由$\eta<\theta$得$\beta\eta\leqslant\alpha$,
记$\varepsilon=\alpha-\beta\eta$, 则有
\[
    \beta\eta+\varepsilon=\alpha<\beta\theta=\beta\eta+\beta\quad\implies\quad
    \varepsilon<\beta.
\]

\noindent{\kaishu 唯一性.}

假设同时存在$\eta,\varepsilon$和$\eta',\varepsilon'$使得
\[\left\{\begin{aligned}
    &\alpha=\beta\eta+\varepsilon\,&&\land\quad\varepsilon<\beta,\\[-3pt]
    &\alpha=\beta\eta'+\varepsilon'\,&&\land\quad\varepsilon'<\beta.
\end{aligned}\right.\]
假设$\eta<\eta'$, 即$\eta^+\leqslant\eta'$, 则
\[
    \alpha=\beta\eta+\varepsilon<\beta\eta+\beta=\beta\eta^+\leqslant
    \beta\eta'\leqslant\beta\eta'+\varepsilon'=\alpha,
\]
矛盾. 同理$\eta'<\eta$也矛盾, 所以$\eta=\eta'$.
进一步显然$\varepsilon=\varepsilon'$. \hfill $\qedsymbol$

\vspace{10pt}

\hypertarget{sec:定理1.5}{}
\noindent\textbf{定理1.5}

任给非零序数$\alpha,\beta$ :
\begin{enumerate}
    \item 如果$\alpha,\beta$都是后继序数, 那么$\alpha\beta$是后继序数,
    \item 如果$\alpha$或$\beta$是极限序数, 那么$\alpha\beta$是极限序数.
\end{enumerate}

\noindent{\kaishu 证明.}

\begin{enumerate}
    \item 如果$\alpha=\gamma^+$, $\beta=\delta^+$, 那么
    \[
        \alpha\beta=\alpha\delta+\alpha=(\alpha\delta+\gamma)^+.
    \]
    \item 如果$\beta$是极限序数, 那么对任意的$\displaystyle
    x<\alpha\beta=\bigcup_{\gamma<\beta}\alpha\gamma$, 存在$\gamma<\beta$使得
    \vspace{3pt}
    \[
        x<\alpha\gamma\quad\implies\quad
        x^+<\alpha\gamma+1\leqslant\alpha\gamma+\alpha=\alpha\gamma^+
        \quad\implies\quad x^+<\alpha\beta,
    \]
    所以$\alpha\beta$是归纳集, 即是极限序数.
    \item 如果$\alpha$是极限序数且$\beta$是后继序数, 即有$\beta=\gamma^+$, 那么
    \[
        \alpha\beta=\alpha\gamma+\alpha
    \]
    是极限序数. \hfill $\qedsymbol$
\end{enumerate}





\newpage

\hypertarget{sec:定理1.6}{}
\noindent\textbf{定理1.6 (序数乘法的右分配律)}

任给序数$\alpha,\beta,\gamma$, 有$\alpha(\beta+\gamma)=\alpha\beta+\alpha\gamma$.

\noindent{\kaishu 证明.}

给定$\alpha,\beta$, 不妨设$\alpha$非零, 对$\gamma$归纳证明.
任取$\gamma$, 假设对任意的$\delta<\gamma$有
$\alpha(\beta+\delta)=\alpha\beta+\alpha\delta$.

\begin{enumerate}
    \item 如果$\gamma=0$, 显然$\alpha(\beta+0)=\alpha\beta=\alpha\beta+\alpha0$.
    
    \vspace{3pt}
    \item 如果$\gamma$是后继序数, 即有$\gamma=\delta^+$, 那么$\delta<\gamma$,
    依假设有
    \[
        \alpha(\beta+\gamma)=\alpha(\beta+\delta)+\alpha=
        \alpha\beta+\alpha\delta+\alpha=\alpha\beta+\alpha\gamma.
    \]

    \vspace{3pt}
    \item 如果$\gamma$是极限序数, 此时$\beta+\gamma$和$\alpha\gamma$也是极限序数.
    
    \vspace{3pt}\hspace{2.17em}
    一方面, 对任意的$\displaystyle\theta<\beta+\gamma=
    \bigcup_{\delta<\gamma}(\beta+\delta)$, 存在$\delta<\gamma$使得
    $\theta<\beta+\delta$, 所以
    \vspace{3pt}
    \[
        \alpha\theta<\alpha(\beta+\delta)=\alpha\beta+\alpha\delta
        \leqslant\sup_{\lambda<\alpha\gamma}(\alpha\beta+\lambda)
        \quad\implies\quad\sup_{\theta<\beta+\gamma}\alpha\theta
        \leqslant\sup_{\lambda<\alpha\gamma}(\alpha\beta+\lambda).
    \]

    \vspace{3pt}\hspace{2.17em}
    另一方面, 对任意的$\displaystyle\lambda<\alpha\gamma=
    \bigcup_{\delta<\gamma}\alpha\delta$, 存在$\delta<\gamma$使得
    $\lambda<\alpha\delta$, 所以
    \vspace{3pt}
    \[
        \alpha\beta+\lambda<\alpha\beta+\alpha\delta=\alpha(\beta+\delta)
        \leqslant\sup_{\theta<\beta+\gamma}\alpha\theta
        \quad\implies\quad\sup_{\lambda<\alpha\gamma}(\alpha\beta+\lambda)
        \leqslant\sup_{\theta<\beta+\gamma}\alpha\theta.
    \]

    \hspace{2.17em}
    综上即得
    \begin{equation*}
        \alpha(\beta+\gamma)=\sup_{\theta<\beta+\gamma}\alpha\theta=
        \sup_{\lambda<\alpha\gamma}(\alpha\beta+\lambda)=
        \alpha\beta+\alpha\gamma.
    \tag*{\qedsymbol}\end{equation*}
\end{enumerate}

\vspace{10pt}

\hypertarget{sec:定理1.7}{}
\noindent\textbf{定理1.7 (序数乘法的结合律)}

任给序数$\alpha,\beta,\gamma$, 有$(\alpha\beta)\gamma=\alpha(\beta\gamma)$.

\noindent{\kaishu 证明.}

给定$\alpha,\beta$, 不妨设$\beta$非零, 对$\gamma$归纳证明.
任取$\gamma$, 假设对任意的$\delta<\gamma$有
$(\alpha\beta)\delta=\alpha(\beta\delta)$.

\begin{enumerate}
    \item 如果$\gamma=0$, 显然$(\alpha\beta)0=0=\alpha(\beta0)$.
    
    \vspace{3pt}
    \item 如果$\gamma$是后继序数, 即有$\gamma=\delta^+$, 那么$\delta<\gamma$,
    依假设有
    \[
        (\alpha\beta)\gamma=(\alpha\beta)\delta+\alpha\beta=
        \alpha(\beta\delta)+\alpha\beta=\alpha(\beta\delta+\beta)=
        \alpha(\beta\gamma).
    \]

    \vspace{3pt}
    \item 如果$\gamma$是极限序数, 此时$\beta\gamma$也是极限序数.
    
    \vspace{3pt}\hspace{2.17em}
    一方面, 对任意的$\delta<\gamma$, 有
    \vspace{3pt}
    \[
        \alpha(\beta\delta)\leqslant\sup_{\lambda<\beta\gamma}\alpha\lambda
        \quad\implies\quad\sup_{\delta<\gamma}\,\alpha(\beta\delta)
        \leqslant\sup_{\lambda<\beta\gamma}\alpha\lambda.
    \]
    
    \vspace{3pt}\hspace{2.17em}
    另一方面, 对任意的$\displaystyle\lambda<\beta\gamma=
    \bigcup_{\delta<\gamma}\beta\delta$, 存在$\delta<\gamma$使得
    $\lambda<\beta\delta$, 所以
    \vspace{3pt}
    \[
        \alpha\lambda<\alpha(\beta\delta)
        \leqslant\sup_{\delta<\gamma}\,\alpha(\beta\delta)
        \quad\implies\quad\sup_{\lambda<\beta\gamma}\alpha\lambda
        \leqslant\sup_{\delta<\gamma}\,\alpha(\beta\delta).
    \]

    \hspace{2.17em}
    综上即得
    \begin{equation*}
        (\alpha\beta)\gamma=\sup_{\delta<\gamma}\,(\alpha\beta)\delta=
        \sup_{\delta<\gamma}\,\alpha(\beta\delta)=
        \sup_{\lambda<\beta\gamma}\alpha\lambda=\alpha(\beta\gamma).
    \tag*{\qedsymbol}\end{equation*}
\end{enumerate}

\end{document}